\chapter{Monotonic Relaxation}
\label{ch:relaxation}

\section{A Terminological Warning}

This chapter's theorem shares a name, and very little else, with
monotonic continuation as developed elsewhere in this program. The two
should not be conflated, and the distinction is stated here before the
theorem itself to prevent exactly that conflation. Monotonic
continuation is an ordering relation, $S_t \le S_{t+1}$, holding between
successive states of a trajectory and asserting that later states
preserve the distinctions from earlier states that remain necessary for
the reachable continuation; it says nothing about any scalar quantity
increasing or decreasing. Monotonic relaxation, the subject of this
chapter, is a claim that a specific scalar functional decreases along a
specific class of trajectories. The two claims are compatible, and one
will be shown below to imply a restricted version of the other under an
additional hypothesis, but they are not restatements of each other and
a trajectory can satisfy either without satisfying the other.

\section{Statement}

\begin{definition}[Closure functional]
For a state $X$, define $\closure{X} = -\log \mu(\reach{X})$, where
$\mu$ is a measure on the state space assigning nonzero mass to any
nonempty reachable set. $\closure{X}$ is large when few futures remain
reachable from $X$ and small when many do; it is a scalar measure of
inadmissibility.
\end{definition}

\begin{definition}[Repair dynamics]
A trajectory follows \emph{repair dynamics} if each transition
$S_i \to S_{i+1}$ is selected, among admissible continuations, to
increase reachable-set volume: $\mu(\reach{S_{i+1}}) \ge
\mu(\reach{S_i})$. This is the operational content of repair within this
program's taxonomy of admissibility-modifying operations: repair is
the restorative case, distinguished from constitutive and regulatory
modification precisely by this volume-increasing character.
\end{definition}

\begin{theorem}[Monotonic Relaxation]
\label{thm:relaxation}
If a trajectory $S_0, \dots, S_t$ follows repair dynamics throughout,
then $\closure{S_i}$ is non-increasing along the trajectory:
$\closure{S_{i+1}} \le \closure{S_i}$ for every $i < t$.
\end{theorem}

\begin{proof}
By definition of repair dynamics, $\mu(\reach{S_{i+1}}) \ge
\mu(\reach{S_i})$ for every $i < t$. Since $-\log$ is a strictly
decreasing function on positive reals, $\mu(\reach{S_{i+1}}) \ge
\mu(\reach{S_i})$ implies $-\log \mu(\reach{S_{i+1}}) \le -\log
\mu(\reach{S_i})$, that is, $\closure{S_{i+1}} \le \closure{S_i}$.
\end{proof}

\section{Discussion}

The proof is immediate once repair dynamics are defined as
volume-increasing, and this is deliberate: the theorem's content is not
in the inference from the definition, which is a one-line consequence of
$-\log$ being decreasing, but in the definition itself being the correct
operational characterization of repair. The definition is not
stipulated arbitrarily; it is inherited directly from the restorative
case of this program's modificatory-operation taxonomy, where repair is
already characterized, independently of this chapter, as the operation
type that restores lost admissible continuations. Theorem
\ref{thm:relaxation} shows that once repair is characterized this way,
monotone decrease of $\closure{\cdot}$ follows automatically; it is not
an additional empirical claim about repair processes but a restatement
of what "restoring lost continuations" already meant, now expressed as
a statement about a scalar functional.

\begin{remark}[Scope of the theorem]
The hypothesis is essential and is not automatically satisfied by an
arbitrary admissible trajectory. Constitutive and regulatory
modifications --- the other two branches of this program's taxonomy of
operations on admissibility structure --- need not increase
$\mu(\reach{\cdot})$, and ordinary state evolution under a fixed $R_t$
can move toward states with strictly smaller reachable sets without
this constituting any failure of admissibility. A claim of the form
"inadmissibility always decreases" would be false; Theorem
\ref{thm:relaxation} is true only of the restorative subclass of
trajectories, and its scope should be stated with this restriction every
time it is invoked.
\end{remark}

\begin{corollary}[Conditional relation to monotonic continuation]
\label{cor:relaxation-continuation}
If a trajectory follows repair dynamics and additionally satisfies
$\reach{S_{i+1}} \supseteq^{*} \reach{S_i}$ in the starred sense used by
monotonic continuation --- that is, $S_{i+1}$ retains the admissible
continuations of $S_i$ that remain significant for the reachable
trajectory, rather than merely having a reachable set of at least equal
measure --- then the trajectory satisfies both $S_i \le S_{i+1}$ in the
sense of monotonic continuation and $\closure{S_{i+1}} \le \closure{S_i}$
in the sense of this chapter. Neither condition implies the other in
general: a trajectory can increase $\mu(\reach{\cdot})$ while failing to
retain the specific significant continuations monotonic continuation
requires, by replacing them with an equal or greater volume of
differently located reachable states; and a trajectory can satisfy
monotonic continuation's retention condition while $\mu(\reach{\cdot})$
decreases, if what is retained is a shrinking but still significant
core.
\end{corollary}

Corollary \ref{cor:relaxation-continuation} is the precise statement of
the relationship gestured at in the opening warning of this chapter.
The two notions of monotonicity coincide only under an additional
hypothesis that must be checked independently in any given application,
and the corollary exists chiefly to prevent the two theorems from being
cited interchangeably.
